\chapter{Gemelos digitales inalámbricos y calibración diferenciable}
\label{ch:wdt}

El capítulo anterior estableció la cadena que conduce desde la propagación hasta la
medición de canal. Sobre esa base, este capítulo introduce el gemelo digital como una
realización computacional de dicha cadena. El interés no reside únicamente en reproducir
el canal observado, sino en conservar la relación entre los parámetros físicos de la escena
y las variaciones de la señal que posteriormente se utilizarán para inferir movimiento y
fisiología.

\section{De la simulación de cobertura al gemelo inalámbrico}

Un simulador de propagación predice $H$ a partir de una escena; un gemelo digital persigue, además, correspondencia con un entorno físico concreto, actualización con mediciones y capacidad de responder a cambios. En el contrato de la \cref{eq:thesis_chain}, no basta con generar mapas plausibles: el modelo debe declarar qué escena, frecuencia, antena y observable reproduce. La formulación mínima es
\begin{equation}
\widehat H_m(f,t)=F_m
\left(\Gamma(t),\vect p_T,\vect p_R,f;
\vect s_R,\vect\theta_E\right),
\end{equation}
donde $\theta_E$ incluye parámetros materiales, dispersión, patrones de antena y otros parámetros físicos no conocidos exactamente.

El reto es que una reconstrucción geométrica visualmente correcta puede producir errores de fase considerables. A frecuencia portadora $f_c$, un error de longitud $\delta L$ produce
\begin{equation}
\delta\phi=-2\pi\frac{\delta L}{\lambda}.
\end{equation}
Por tanto, los errores centimétricos o incluso milimétricos pueden ser importantes para el sensado de movimientos pequeños. Esta observación motiva dos líneas complementarias: la calibración diferenciable de propiedades electromagnéticas y la calibración que considera errores locales de fase.

\section{Calibración como problema inverso}

Sea $\vect\theta$ el vector de parámetros del gemelo y $\mathcal D_{\rm cal}$ un conjunto de mediciones alineadas. Se define
\begin{equation}
\vect\theta^\star=\arg\min_{\vect\theta}\;
\sum_{i\in\mathcal D_{\rm cal}}
\calL\left(
H_{m,i}^{\rm real},
F_m(\Gamma_i,\vect p_{T,i},\vect p_{R,i},f;
\vect s_{R,i},\vect\theta)
\right)+\beta R(\vect\theta).
\label{eq:calibration_inverse}
\end{equation}
Para una pérdida real se vectorizan parte real e imaginaria, $\vect u=[\Re H;\Im H]$, y la diferenciabilidad permite obtener
\begin{equation}
\nabla_{\vect\theta}\calL
=\left(\frac{\partial\vect u}{\partial\vect\theta^T}\right)^T
\nabla_{\vect u}\calL.
\end{equation}
Trabajos recientes de calibración con \ac{DRT} muestran que las propiedades materiales, patrones de antena y parámetros de dispersión pueden optimizarse utilizando canales reales \citep{Hoydis2024DiffRT}; otros introducen errores locales de fase para modelar discrepancias geométricas pequeñas que una calibración de potencia no puede absorber \citep{Ruah2024Phase}. En esta tesis se consideran ambas familias como componentes consecutivos, no sustitutos.

\section{Separación de fuentes de error}

Conviene escribir
\begin{equation}
H_{\rm real}-H_{\rm gemelo}=
\Delta H_{\rm geom}+
\Delta H_{\rm mat}+
\Delta H_{\rm scat}+
\Delta H_{\rm ant}+
\Delta H_{\rm hw}+
\epsilon_{\rm model}.
\label{eq:domain_gap_decomp}
\end{equation}
Esta suma es una contabilidad conceptual o una linealización local, no una descomposición identificable a partir de $H$ por sí solo. Un modelo flexible todavía puede usar materiales para compensar coordenadas. Por ello la calibración será jerárquica: registro y sincronización; antenas y escala; materiales y scattering; y sólo al final residuales de fase/modelo. Priors, mediciones auxiliares e intervenciones separadas anclarán cada bloque.

\section{Qué resuelve la calibración diferenciable}
\label{sec:literature_evidence_hierarchy}

Sionna RT proporciona un trazador de rayos diferenciable y una interfaz entre escena,
antenas y canal \citep{Hoydis2023SionnaRT}. Sobre esa base, la calibración de Hoydis
y colaboradores trata materiales y dispersión como parámetros optimizables frente a
mediciones \citep{Hoydis2024DiffRT}. Su contribución metodológica es hacer explícito
el modelo directo y propagar gradientes a parámetros de la escena. El observable,
la función de pérdida y la regularización determinan, sin embargo, qué se recupera:
un mínimo de potencia no identifica por sí solo permitividad, rugosidad y error de
registro.

Ruah y colaboradores introducen variables locales de error de fase para absorber
discrepancias geométricas pequeñas \citep{Ruah2024Phase}. Esta familia responde a
un problema que la calibración exclusiva de materiales no puede resolver con
facilidad: errores sublongitud de onda producen rotaciones grandes aunque potencia
y retardos gruesos sean razonables. La ventaja tiene un costo de identificabilidad.
Una corrección flexible puede reunir desplazamiento de interacción, error de reloj
y mecanismos ausentes; debe probarse fuera de las observaciones que la estimaron.

La diferencia entre las tres estrategias consideradas en ese marco se entiende al
introducir un error circular $\varepsilon_p$ en cada contribución simulada,
\begin{equation}
H^{\rm mod}(f)=\sum_p \alpha_p(f)
\exp\{-j2\pi f\tau_p+j\varepsilon_p\}.
\label{eq:path_phase_error_model}
\end{equation}
La calibración que ignora el error de fase (PEOC) fija
$\varepsilon_p=0$ y obliga a los materiales a explicar también cualquier error
geométrico. UPEC marginaliza fases sin estructura mediante una distribución
uniforme. PEAC conserva información parcial con una distribución de von Mises,
\begin{equation}
p(\varepsilon_p\mid\mu_p,\kappa_p)=
\frac{\exp\{\kappa_p\cos(\varepsilon_p-\mu_p)\}}
{2\pi I_0(\kappa_p)},
\label{eq:von_mises_path_error}
\end{equation}
donde $\kappa_p=0$ recupera el caso uniforme y una concentración grande expresa
confianza alrededor de $\mu_p$. Esta jerarquía aclara el compromiso: PEOC puede
sesgar materiales cuando la fase RT es incorrecta; UPEC evita esa confianza, pero
descarta estructura; PEAC puede aprovecharla, aunque introduce variables latentes,
dependencia de inicialización y posibles óptimos locales. Sin información adicional,
una fase por trayectoria puede ser numéricamente estimable y aun carecer de una
interpretación física única.

Estas obras motivan una calibración escalonada. Primero se fijan convención,
registro y respuesta de cadena; después se ajustan parámetros globales o compartidos
de materiales y antenas; finalmente se estudia discrepancia local con priors y
particiones espaciales. La evidencia necesaria aumenta en el mismo orden:
correspondencia macroscópica, generalización estática y, para esta tesis, respuesta
diferencial ante una intervención. Que $H_{\rm model}$ sea preciso en los puntos de
calibración no implica que su derivada respecto de $q$ sea correcta.

\subsection{Escalera de capacidad B0--B2}

La evidencia experimental utilizará tres configuraciones anidadas inspiradas en
la calibración de Hoydis y colaboradores y en su implementación reproducible
\citep{Hoydis2024DiffRT,NVIDIA2023DiffRTCalibration}. Conviene introducirlas antes
de presentar sus resultados. Las tres conservan la escena, posiciones, soporte de
trayectorias, frecuencia y operador de síntesis. La diferencia reside en la
representación electromagnética evaluada en los puntos de interacción.

En B0, las superficies conservan propiedades nominales de la recomendación ITU;
no existe una red neuronal ni se aprende variación espacial de materiales. Un
factor escalar global iguala la escala de potencia entre medición y simulación,
pero no altera retardos, fases relativas ni materiales. B0 es, por tanto, la
referencia determinista nominal.

En B1, cada objeto de la escena recibe un vector constante de parámetros
electromagnéticos. En la configuración empleada, las variables físicamente activas
son permitividad relativa real $\epsilon_r$ y conductividad $\sigma$; permanecen
constantes sobre todos los puntos de un mismo objeto y se optimizan por gradiente a
través del trazador diferenciable. B1 es una calibración paramétrica por objeto,
no una red que consulte coordenadas.

En B2, la parametrización constante se sustituye por un campo aprendido
\begin{equation}
(\epsilon_r(\vect x),\sigma(\vect x),s(\vect x),\chi(\vect x))
=g_{\vect\psi}\!\left(\gamma(\vect x)\right),
\label{eq:neural_material_field}
\end{equation}
donde $\vect x$ es el punto de interacción, $\gamma$ es una codificación posicional
senoidal y $g_{\vect\psi}$ es un \ac{MLP}. B2 sigue usando trayectorias explícitas:
la red parametriza materiales locales, no reemplaza el trazador por un campo
neuronal de canal. La arquitectura, las salidas activas y el protocolo de
optimización se especifican junto con los resultados en el \cref{ch:avances}.

La numeración comienza en cero: B2 es la tercera configuración. No existe una B3
en la serie experimental documentada. Esta aclaración impide confundir B2 con
NeRF$^2$, Learnable WDT u otras arquitecturas que se discuten después como
comparaciones pendientes.

\section{Gemelos aprendidos y alcance de la comparación}

Los campos implícitos, gemelos aprendibles y representaciones centradas en objetos
amplían el conjunto de funciones con las que puede aproximarse un entorno
\citep{Zhao2023NeRF2,Jiang2025LearnableWDT,Chen2025RFScape,Zhang2026RadTwin}.
Resuelven problemas relacionados, pero no equivalentes: interpolación del campo,
reconstrucción electromagnética, edición de objetos y transferencia entre escenas.
Su comparación requiere igualar entradas disponibles, posiciones de entrenamiento,
observable y partición. El \cref{ch:latent_env} desarrolla estas familias sin
presentarlas como sustitutos directos de la calibración física.

La literatura de simulación dinámica e inversión
\citep{Zheng2026WaveVerse,Han2026RayLoc,Chen2026RFDT,Chen2025InverTwin}
demuestra que un modelo directo estructurado puede apoyar tareas posteriores.
Para el argumento de esta tesis, esos resultados establecen plausibilidad
metodológica, no evidencia de que el Jacobiano del banco propio sea correcto.
El reflector y el fantoma se introducen precisamente para cerrar esa brecha.

\section{DICHASUS dc01 como banco de pruebas del gemelo estático}

La primera validación desarrollada utiliza DICHASUS dc01 y la escena correspondiente empleada en trabajos de calibración diferenciable \citep{Euchner2023DICHASUSdcxx,Hoydis2024DiffRT}. El banco reportado comprende $N=10000$ posiciones y conserva el objeto experimental
\begin{equation}
\mathcal D_{\rm dc01}
=\left\{
\vect p_i,\;H_{m,i}^{\rm real},\;\calP_{m,i}^{\rm RT},\;H_{m,i}^{\rm RT}
\right\}_{i=1}^{N}.
\label{eq:dc01_dataset}
\end{equation}
El soporte geométrico y el conjunto de trayectorias radio--dependiente se conservan como interfaz física, en lugar de almacenar únicamente el canal sintetizado. Esto permite modificar posteriormente materiales, frecuencia, observación y dinámica sin confundir el mundo común $\Gamma$ con una realización particular $\calP_m$.

La comparación B0--B2 y los diagnósticos de fase se documentan en el \cref{ch:avances}. DICHASUS es un banco de sondeo OFDM: su respuesta en frecuencia resulta adecuada para estudiar la propagación, mientras la validación del operador de observación de dispositivos Wi-Fi comerciales requiere pruebas específicas. Conservar la geometría tampoco autoriza reutilizar sin cambios las ganancias y trayectorias efectivas al pasar a LoRa.

\section{Fidelidad de campo y fidelidad diferencial}

Para comunicaciones suele ser suficiente comparar $H_{\rm sim}$ y $H_{\rm real}$. Para el sensado se requiere una condición adicional. Si $q$ es una variable física de interés,
\begin{equation}
J_q=\frac{\partial H}{\partial q},
\end{equation}
y se propone medir
\begin{equation}
E_J=
\frac{\|J_q^{\rm model}-J_q^{\rm ref}\|_F}
{\|J_q^{\rm ref}\|_F+\epsilon}.
\label{eq:jacobian_error}
\end{equation}
Un modelo con bajo \ac{NMSE} de canal, pero con derivadas incorrectas respecto del movimiento humano, puede ser inadecuado para el sensado. Por ello, la tesis distingue entre \emph{fidelidad del modelo directo} y \emph{fidelidad diferencial}.

$J_q^{\rm ref}$ se estimará con perturbaciones mecánicas bidireccionales, paso declarado, repeticiones y alineamiento de fase/ganancia definido antes de observar el resultado. Como el error relativo de la \cref{eq:jacobian_error} es inestable cuando $\|J_q^{\rm ref}\|$ es cercano a cero, se reportará también error absoluto con unidades, correlación compleja y error ponderado por la tarea. El término $\epsilon$ se fijará en las mismas unidades que la norma de referencia y no como constante adimensional arbitraria.

\section{Reconstrucción métrica y registro del gemelo}
\label{sec:geometric_metrology}

La geometría debe representarse con unidades, sistema de coordenadas e incertidumbre
verificables. Se distinguen tres objetivos de diseño: escala global de paredes,
puertas y mobiliario del orden de 1--2 cm; registro de antenas del orden de 2--5 mm;
y referencia local de desplazamiento del orden de 0.1--1 mm. Son rangos de trabajo
que deberán demostrarse mediante controles independientes, no especificaciones
atribuidas a cualquier LiDAR, cámara o actuador. Un experimento con pasos de
0.25 mm sólo es interpretable si su incertidumbre efectiva permite resolverlos.

Se fijará un marco global $W$ y se almacenará, para cada sensor $s$,
\begin{equation}
T_{W\leftarrow s}=\begin{bmatrix}R_{Ws}&\mathbf t_{Ws}\\0&1\end{bmatrix}\in SE(3),
\qquad \widetilde{\mathbf p}_W=T_{W\leftarrow s}\widetilde{\mathbf p}_s,
\label{eq:scene_extrinsics}
\end{equation}
con $R_{Ws}\in SO(3)$. Las transformaciones compuestas deberán respetar esa dirección.
Se registrarán extrínsecos de cámara, LiDAR, radar, transmisores, receptores y antenas;
el punto de montaje no se confundirá con el centro de fase, que puede depender de
frecuencia y dirección. Marcadores rígidos, distancias de control y mediciones repetidas
permitirán verificar el registro. Un residual pequeño de ICP mide concordancia entre
nubes, pero no garantiza exactitud absoluta. Se reservarán, cuando el recinto lo permita,
10--20 distancias de control que no se hayan utilizado para ajustar la reconstrucción.

La comparación geométrica avanzará desde primitivas medidas hasta RGB-D, LiDAR
regularizado para RF y calibración electromagnética, siempre en el mismo marco global.
Como el último nivel modifica geometría y materiales, se requiere una ablación
factorial que separe registro, regularización y ajuste electromagnético. Los puntos
de control reservados medirán generalización y no se utilizarán para corregir la
escena. El \cref{ch:platform} especifica la realización metrológica de estos niveles.

RFDT-Channel ofrece una referencia reciente para construir escenas RF a partir de
RGB y LiDAR y simular canales interiores a 28 GHz \citep{Yao2026RFDTChannel}.
El gemelo de Wang y colaboradores \citep{Wang2026WiFiTwin} contrasta predicción de señal con RSSI real.
Ambos orientan el flujo geométrico, pero una validación de potencia o un canal simulado
no establecen exactitud de fase medida ni fidelidad de micromovimiento de este banco.
La comparación G0--G4 se evaluará sucesivamente en error geométrico, potencia,
retardos, canal relativo y respuesta diferencial. No se exigirá resolver primero
la fase absoluta de toda la habitación para iniciar una prueba local controlada.

Esta distinción cierra el paso entre el modelado electromagnético y la representación: una
aproximación útil debe conservar no sólo valores de canal, sino también las variaciones
que contienen información acerca del estado humano. El capítulo siguiente examina qué
estructuras explícitas, implícitas o híbridas pueden satisfacer ese requisito.
